\section{Motivation}
\label{sec:motivation}

The core motivation of this work is both empirical and methodological:
empirically, because blockchain fraud continues to inflict large-scale
financial harm that existing single-modal detectors cannot fully address;
methodologically, because the research community lacks an openly seeded,
multi-modal baseline against which future detectors can be honestly
compared.

\subsection{The Multi-Modal Nature of On-Chain Fraud}

Ethereum fraud is not a single phenomenon.  It manifests in at least two
distinct and complementary signatures.

\emph{Behavioural signatures} emerge in the transaction graph.
A phishing contract receives funds from hundreds of victims before routing
them to a handful of cashout addresses---a star-shaped in-flow that is
anomalous regardless of what the contract code says.
A rug-pull token accretes liquidity from many investors and then drains it
to a single deployer wallet in one block.
These patterns are detectable purely from the structure of who sent what
to whom.

\emph{Semantic signatures} emerge in the contract source code.
A flash-loan exploit embeds a reentrancy callback that steals funds
before a balance update---a logic vulnerability invisible in the on-chain
footprint of a single malicious transaction, yet unambiguous to a code
analyser.
A honeypot contract contains a hidden owner function that prevents
withdrawals from non-deployer addresses; its transaction pattern may
appear entirely normal.

The key insight is that these two signature types are complementary, not
redundant.
A contract that mimics the transaction behaviour of a legitimate exchange
while embedding a backdoor in its source code cannot be caught by a
graph-only detector.
A contract whose source code is a clean fork of a legitimate template but
whose transaction neighbourhood includes known fraud addresses cannot be
caught by a code-only detector.
Only a detector with access to both modalities simultaneously can
disambiguate such cases---a claim our ablation study validates
empirically (Section~\ref{subsec:ablation_module}).

\subsection{The Data Pipeline Determines the Ceiling}

The quality of a fraud detector is bounded by the quality of its training
data.  Two pipeline decisions have an outsized effect.

First, \emph{label provenance}.
Using attacker-wallet addresses as fraud seeds---as is common in Forta-based
pipelines---produces addresses with no verified source code, because
attackers rarely publish their tools on Etherscan.  The fusion model
cannot be trained on such samples.  We instead seed from
\emph{victim contracts}: real protocol contracts that were exploited, each
with verified Solidity source and confirmed on-chain transaction history.
This ``victim-first'' strategy is what makes multi-modal training
tractable on real Ethereum data.

Second, \emph{collection strategy}.
Random sampling across the Ethereum blockchain produces a dataset in which
fewer than 0.002\% of addresses are labelled as fraud.
Training on such data teaches a classifier to predict benign for
everything.  Our needle-first approach---retrieving the transaction
neighbourhood of known fraud and benign seed contracts---produces a dense,
label-rich subgraph in which every transaction is connected to at least
one seed address.

\subsection{Reproducibility as a Scientific Requirement}

Blockchain fraud detection is an applied field with real financial stakes.
A model that exists only as an entangled notebook cannot be audited by a
regulator, extended by a competitor, or retrained against new exploit
signatures by a practitioner.

We therefore treat reproducibility as a first-class design constraint.
The two seed arrays are explicit Python lists committed to the repository.
Every pipeline step is a numbered script that reads from deterministic
inputs and writes to deterministic outputs.
The train/val/test split is seeded and serialised before any model
training occurs.
These design choices are not merely hygienic---they make the experimental
claims of this paper verifiable.

The three motivations above are tightly coupled: the right architecture
can only be trained on the right data, and its scientific value depends
on reproducibility.  The remainder of the paper develops each dimension.
